\section{Conclusion}
\label{sec:conclusion}

This study examines when operational information is sufficient for decisions in rolling airline disruption recovery. Using a common multidimensional consequence criterion and two distinct decision tasks---fixed-budget recovery screening and effective local-recovery targeting---the empirical design evaluates delay-based priorities, temporal information, and probabilistic state representations across U.S. airline operations in 2017--2022. The central question is not whether a representation retains more operating-state detail, but whether that detail changes the allocation of limited managerial attention or the selection of a worthwhile local-recovery target.

The empirical conclusions will be organized around three related findings. First, the six-year screening comparisons establish \emph{[E1 finding pending: when delay-based screening does or does not preserve consequence-based shortlists and their reference consequence coverage]}. Second, the full eligible recovery cohorts show \emph{[E2 finding pending: the observed relationship between baseline consequence priority and remaining recoverable opportunity]}; these are distinct managerial criteria rather than interchangeable measures. Third, the temporal and distributional comparisons, together with objective-specific and operating-regime analyses, indicate \emph{[E2--E4 finding pending: which information distinctions alter decisions, under which objectives and operating conditions]}. These findings must be interpreted through the corresponding decision alternatives and reference losses, rather than inferred from predictive accuracy or representational richness alone.

The managerial implication is \emph{decision-matched information provision}: the analytical detail supplied to an airline operations control center should correspond to the decision it supports. Initial screening, detailed consequence assessment, and local-recovery targeting need not require the same information depth at every operating epoch. Translating this principle into executable airline recovery would require carrier-internal operating information, resource-explicit recovery actions, and evaluation against realized operational outcomes.

% SCIENTIFIC COMPLETION GATE: Replace all three visible bracketed empirical
% placeholders only after the six-year E1--E4 results and Marginal/support
% contracts have been audited. Do not transfer 2019-only findings to 2017--2022.
% The core conceptual claims and concluding managerial principle do not imply
% measured latency savings, realized operational benefits, or causal regime effects.
